% Môn: Vật lý
% Lớp: 11
% Chương: Chương I. Dao động
% Bài: L11C1 Bài 1. Dao động điều hòa · Xác định các đại lượng dựa vào đồ thị dao động ($x-t$, $v-t$, $a-t$).
% ===== Câu 25 =====
% ID: L11C1B1-176-DS
% Mức: TH
\dangbt{Xác định các đại lượng dựa vào đồ thị dao động ($x-t$, $v-t$, $a-t$).}
\begin{ex}
% ID: L11C1B1-176-DS
% Mức: TH
Đồ thị $x-t$ cho biên độ $8\,\text{cm}$, chu kì $1\,\text{s}$ và ban đầu ở biên dương.
\choiceTF
{\True Biên độ $8\,\text{cm}$.}
{\True Tần số $1\,\text{Hz}$.}
{\True Vận tốc ban đầu bằng 0.}
{\True Có thể viết $x=8\cos(2\pi t/1)$ cm.}
\end{ex}

% ===== Câu 26 =====
% ID: L11C1B1-177-DS
% Mức: TH
\begin{ex}
% ID: L11C1B1-177-DS
% Mức: TH
Đồ thị li độ–thời gian cho biết $A=4$ cm, hai đỉnh dương liên tiếp cách nhau $1\,\text{s}$ và tại $t=0$ vật ở biên dương. Xác định $T,f,\omega$ và viết phương trình dao động.
\choiceTF

\end{ex}

% ===== Câu 29 =====
% ID: L11C1B1-178-DS
% Mức: VD
\begin{ex}
% ID: L11C1B1-178-DS
% Mức: VD
Đồ thị $x-t$ cho biên độ $9\,\text{cm}$, chu kì $2\,\text{s}$ và ban đầu ở biên dương.
\choiceTF
{\True Biên độ $9\,\text{cm}$.}
{\True Tần số $0.5\,\text{Hz}$.}
{\True Vận tốc ban đầu bằng 0.}
{\True Có thể viết $x=9\cos(2\pi t/2)$ cm.}
\end{ex}

% ===== Câu 30 =====
% ID: L11C1B1-179-DS
% Mức: VD
\begin{ex}
% ID: L11C1B1-179-DS
% Mức: VD
Đồ thị li độ–thời gian cho biết $A=6$ cm, hai đỉnh dương liên tiếp cách nhau $2\,\text{s}$ và tại $t=0$ vật ở biên dương. Xác định $T,f,\omega$ và viết phương trình dao động.
\choiceTF

\end{ex}

% ===== Câu 33 =====
% ID: L11C1B1-180-DS
% Mức: VD
\begin{ex}
% ID: L11C1B1-180-DS
% Mức: VD
Một vật dao động điều hòa trên trục Ox với hình vẽ
\choiceTF

\loigiai{
A. Sai — Tại $t = 0$, $x = 5\cos\!\left(-\dfrac{\pi}{3}\right) = 2,5$ cm, không phải $-2,5$ cm.
B. Sai — Tại $t = 1$ s, $x = 5\cos\!\left(\pi \cdot 1 - \dfrac{\pi}{3}\right) = 5\cos\!\left(\dfrac{2\pi}{3}\right) = -2,5$ cm, không phải $2,5$ cm.
C. Sai — Từ $t = 0$ đến $t = 1,5$ s, pha tăng từ $-\dfrac{\pi}{3}$ đến $\dfrac{7\pi}{6}$, vật đổi chiều chuyển động 2 lần, không phải 4 lần.
D. Đúng — Tại $t = \dfrac{11}{6}$ s, $x = 5\cos\!\left(\pi \cdot \dfrac{11}{6} - \dfrac{\pi}{3}\right) = 5\cos\!\left(\dfrac{3\pi}{2}\right) = 0$, nhưng đây là lần li độ bằng 0 thứ hai, không phải lần thứ nhất.
}
\end{ex}

% ===== Câu 46 =====
% ID: L11C1B1-236-TLN
% Mức: TH
\begin{ex}
% ID: L11C1B1-236-TLN
% Mức: TH
Trên đồ thị li độ–thời gian, hai đỉnh dương liên tiếp ở $t=1$ s và $t=2$ s. Chu kì theo đơn vị s bằng bao nhiêu?
\shortans{31/12/1899}
\end{ex}

% ===== Câu 48 =====
% ID: L11C1B1-237-TLN
% Mức: VD
\begin{ex}
% ID: L11C1B1-237-TLN
% Mức: VD
Trên đồ thị li độ–thời gian, hai đỉnh dương liên tiếp ở $t=1$ s và $t=3$ s. Chu kì theo đơn vị s bằng bao nhiêu?
\shortans{2}
\end{ex}

% ===== Câu 50 =====
% ID: L11C1B1-238-TLN
% Mức: VD
\begin{ex}
% ID: L11C1B1-238-TLN
% Mức: VD
Trên đồ thị li độ–thời gian, hai đỉnh dương liên tiếp ở $t=1$ s và $t=5$ s. Chu kì theo đơn vị s bằng bao nhiêu?
\shortans{4}
\end{ex}

% ===== Câu 59 =====
% ID: L11C1B1-239-TN
% Mức: TH
\begin{ex}
% ID: L11C1B1-239-TN
% Mức: TH
Từ đồ thị dao động, khoảng thời gian giữa hai đỉnh dương liên tiếp là $1\,\text{s}$ và li độ cực đại là $4\,\text{cm}$. Chu kì bằng
\choice
{\True $1\,\text{s}$}
{$0.5\,\text{s}$}
{$2\,\text{s}$}
{$4\,\text{s}$}
\end{ex}

% ===== Câu 61 =====
% ID: L11C1B1-240-TN
% Mức: VD
\begin{ex}
% ID: L11C1B1-240-TN
% Mức: VD
Từ đồ thị dao động, khoảng thời gian giữa hai đỉnh dương liên tiếp là $2\,\text{s}$ và li độ cực đại là $6\,\text{cm}$. Chu kì bằng
\choice
{\True $2\,\text{s}$}
{$1\,\text{s}$}
{$4\,\text{s}$}
{$8\,\text{s}$}
\end{ex}

% ===== Câu 63 =====
% ID: L11C1B1-241-TN
% Mức: VD
\begin{ex}
% ID: L11C1B1-241-TN
% Mức: VD
Từ đồ thị dao động, khoảng thời gian giữa hai đỉnh dương liên tiếp là $4\,\text{s}$ và li độ cực đại là $8\,\text{cm}$. Chu kì bằng
\choice
{\True $4\,\text{s}$}
{$2\,\text{s}$}
{$8\,\text{s}$}
{$16\,\text{s}$}
\end{ex}
